\documentclass[10pt,twocolumn]{article}

\usepackage[utf8]{inputenc}
\usepackage[T1]{fontenc}
\usepackage{amsmath,amssymb,amsfonts,amsthm}
\usepackage{graphicx}
\usepackage{booktabs}
\usepackage{multirow}
\usepackage{geometry}
\usepackage{hyperref}
\usepackage{cite}
\usepackage{caption}

\geometry{a4paper, margin=0.70in, columnsep=0.26in}

\hypersetup{
    colorlinks=true,
    linkcolor=blue,
    citecolor=blue,
    urlcolor=blue
}

\newtheorem{theorem}{Theorem}
\newtheorem{proposition}{Proposition}
\newtheorem{lemma}{Lemma}
\newtheorem{definition}{Definition}

\title{\textbf{Regional Simplicial Decomposition and Post-Inversion Empirical Bayes Calibration for Zero-Shot and Few-Shot Topological Neurodynamics in Multichannel EEG}}

\author{
    \textbf{Ronen Weinberg} \\
    Independent Researcher, SYMETR Research Lab, Jerusalem, Israel \\
    Repository: \texttt{https://github.com/RW-101/TRR3-formalization-}
}

\date{October 2026}

\begin{document}

\maketitle

\begin{abstract}
Zero-lag volume conduction and inter-subject baseline shift remain two primary mathematical obstacles to deploying algebraic topological data analysis (TDA) in real-time brain-computer interfaces (BCIs). While Vietoris--Rips filtration over the finite field $\mathbb{F}_2$ induced by Imaginary Coherency ($\mathrm{ImCoh}$) guarantees exact invariance to instantaneous volume conduction ($\Delta\phi = 0$), scaling global 20-channel minimum spanning tree (MST) persistence $L_{\mathrm{tot}}(H_0)$ from small homogeneous benchmarks ($N=5$, $\mathrm{AUC}=0.764$) to heterogeneous cohorts ($N=20$, $1{,}160$ epochs) causes severe accuracy degradation ($\mathrm{AUC}=0.650$) and sign inversions in $20\%$ of subjects ($4/20$). In this paper, we prove that this degradation stems from three mechanisms: (i) zero-lag denominator inflation under high-amplitude standing waves, (ii) sub-epoch phase cancellation across 4-second windows, and (iii) combinatorial spatial MST dilution across anterior artifact channels. We introduce the \textbf{TRR 3.0 Stage 7} architecture to resolve all three bottlenecks. First, we formulate Sub-Window Rectified and Lag-Orthogonalized Imaginary Coherency ($\mathrm{rLagImCoh}$), proving that it preserves exact zero-lag isolation ($\mathrm{rLagImCoh}(\Delta\phi=0) \equiv 0$) while preventing standing-wave denominator suppression. Second, we introduce an anatomical simplicial decomposition partitioning the 10--20 montage into an 8-channel occipito-parietal complex $\mathcal{K}_{\mathrm{post}}$ and a 12-channel anterior reference complex $\mathcal{K}_{\mathrm{ant}}$, reducing 2-simplex boundary complexity by $4.13\times$ ($1{,}140 \to 276$ triangles). Third, we derive a strictly causal Post-Inversion Empirical Bayes Stein Shrinkage Fusion protocol under Leave-One-Subject-Out Cross-Validation (LOSO-CV) that prevents covariance-inversion gate cancellation. Audited across $N=20$ subjects ($1{,}160$ epochs) from the PhysioNet EEGBCI database, the pure $\mathbb{F}_2$ topological spatial gradient elevates zero-shot ($0\text{ s}$) accuracy from $0.650$ to $0.837$ (mean within-subject $\mathrm{AUC} = 0.927 \pm 0.022$) and eliminates all subject polarity inversions ($20/20$ consistent directionality, paired Wilcoxon $p = 1.91 \times 10^{-6}$). Under causal prefix calibration, Pure Topological TRR ($\mathbb{F}_2$ homology only, $0\%$ local power) scales monotonically in pooled accuracy to $0.874$ at $22\text{ s}$ ($k=10$) and $0.895$ at $60\text{ s}$, while the Hybrid Spectral-Topological TRR architecture achieves $\mathrm{AUC} = 0.869$ at zero-shot ($0\text{ s}$), $0.897$ at $8\text{ s}$ ($k=3$), $0.928$ at $22\text{ s}$ ($k=10$), and $0.954$ at $60\text{ s}$ ($d = +2.34$).
\end{abstract}

\textbf{Keywords:} Relational Resolution Theory (TRR 3.0), Persistent Homology over $\mathbb{F}_2$, Imaginary Coherency, Volume Conduction Invariance, Empirical Bayes Stein Shrinkage, Zero-Shot BCI Calibration, Multichannel EEG.

---

\section{Introduction}

Topological Data Analysis (TDA) provides a coordinate-free mathematical framework for extracting multiscale relational invariants from point clouds and weighted networks. However, generic TDA is agnostic to the biophysics of electromagnetic field propagation in neural tissue. **Relational Resolution Theory 3.0 (TRR 3.0)**~\cite{weinberg2026trr1} bridges this gap by embedding exact simplicial homology over the prime field $\mathbb{F}_2 = \{0,1\}$ within a physically constrained phase-lagged metric space designed specifically for multichannel electroencephalography (EEG). By constructing the Vietoris--Rips filtration over Imaginary Coherency ($\mathrm{ImCoh}$)~\cite{nolte2004}, TRR 3.0 enforces an exact physical boundary condition: any instantaneous linear mixing of cortical dipoles ($\Delta\phi = 0 \pmod\pi$) yields zero imaginary cross-spectrum and maps to the metric ceiling $D(u,v) = 1.0$, guaranteeing immunity against volume-conduction short-circuiting.

Despite this formal guarantee, scaling head-wide topological invariants across diverse human populations reveals a critical bottleneck. On a 5-subject benchmark (\texttt{S001--S005}, $290$ epochs) from the PhysioNet EEG Motor Movement/Imagery Database (EEGBCI)~\cite{schalk2004}, uncalibrated 20-channel zero-dimensional persistence $L_{\mathrm{tot}}(H_0)$ separates resting-state eyes-open ($\mathrm{EO}$) from eyes-closed ($\mathrm{EC}$) neurodynamics with $\mathrm{ROC\text{-}AUC} = 0.764$, rising to $0.788$ after 60 seconds of Empirical Bayes calibration. However, when the cohort is expanded fourfold to $N=20$ subjects (\texttt{S001--S020}, $1{,}160$ epochs), global uncalibrated accuracy drops to $\mathrm{AUC} = 0.650$ ($d = -0.60$), and 60-second univariate calibration reaches only $0.654$.

In this study, we prove mathematically and demonstrate empirically that this degradation is not a limitation of $\mathbb{F}_2$ persistent homology, but rather a consequence of three pre-filtration and spatial aggregation bottlenecks. We then present the complete **TRR 3.0 Stage 7** mathematical architecture that resolves these bottlenecks and achieves state-of-the-art zero-shot and few-shot classification across all $20/20$ subjects.

---

\section{Mathematical Diagnosis of Global Degradation on Heterogeneous Cohorts}

Let $\mathbf{x}(t) \in \mathbb{R}^{N_{\mathrm{ch}}}$ ($N_{\mathrm{ch}}=20$) denote the average-referenced, bandpass-filtered EEG time series over a 4-second epoch of $W = 640$ samples ($f_s = 160\text{ Hz}$), and let $z_u(t) = x_u(t) + i\,\mathcal{H}[x_u](t)$ denote the analytic signal obtained via the discrete Hilbert transform $\mathcal{H}$.

\subsection{Zero-Lag Denominator Inflation}
In standard 4-second Imaginary Coherency~\cite{nolte2004}, the functional coupling between channels $u$ and $v$ is evaluated as:
\begin{equation}
    \mathrm{ImCoh}_{uv} = \frac{|\operatorname{Im}(S_{uv})|}{\sqrt{P_u P_v}}, \qquad S_{uv} = \frac{1}{W}\sum_{t=1}^{W} z_u(t)z_v^*(t),
    \label{eq:old_imcoh}
\end{equation}
where $P_u = S_{uu} = \operatorname{Re}(S_{uu})$ is the total narrowband power. During eyes-closed ($\mathrm{EC}$) alpha synchronization ($8\text{--}13\text{ Hz}$), the analytic signal at electrodes $u,v$ is a superposition of a dominant zero-lag standing-wave projection $A_0(t)$ ($\Delta\phi = 0$, volume-conducted from radial occipito-parietal dipoles) and true phase-lagged travelling cortical waves $\epsilon_u(t), \epsilon_v(t)$ ($\Delta\phi \neq 0$):
\begin{equation}
    z_u(t) = A_0(t) + \epsilon_u(t), \qquad z_v(t) = A_0(t) + \epsilon_v(t).
\end{equation}
Because $\operatorname{Im}(A_0(t)A_0^*(t)) \equiv 0$, the numerator $|\operatorname{Im}(S_{uv})|$ scales exclusively with the lagged interaction $|\operatorname{Im}(\langle \epsilon_u \epsilon_v^* \rangle)|$. In contrast, the denominator $\sqrt{P_u P_v} \approx |A_0|^2 + |\epsilon|^2$ scales linearly with the zero-lag standing-wave power $|A_0|^2$. In subjects whose Berger reaction is dominated by a high-amplitude radial standing wave, $|A_0|^2$ increases by a factor of $5\text{--}10\times$ during $\mathrm{EC}$, inflating the denominator of Eq.~\eqref{eq:old_imcoh} faster than the numerator. Consequently, $\mathrm{ImCoh}_{uv}$ paradoxically decreases during $\mathrm{EC}$, and the metric distance $D(u,v) = 1 - \mathrm{ImCoh}_{uv}$ shifts in the inverted direction ($d > 0$ rather than $d < 0$). As documented in Table~\ref{tab:diag}, exactly $4/20$ subjects (\texttt{S006}, \texttt{S007}, \texttt{S012}, \texttt{S015}) exhibit inverted effect signs ($d > 0$) under global 4-second $\mathrm{ImCoh}$.

\subsection{Sub-Epoch Phase Cancellation}
Cortical phase-lead relationships are metastable, reversing direction every $300\text{--}800\text{ ms}$ ($\sin\Delta\phi(t) > 0 \leftrightarrow \sin\Delta\phi(t) < 0$). Linearly averaging $z_u(t)z_v^*(t)$ over a long $W = 4\text{ s}$ window ($\approx 40$ alpha cycles) \textit{before} taking the modulus $|\operatorname{Im}(\cdot)|$ results in destructive algebraic cancellation between positive and negative phase-lag episodes.

\subsection{Combinatorial Spatial MST Dilution}
Zero-dimensional total persistence $L_{\mathrm{tot}}(H_0)$ over $N_{\mathrm{ch}} = 20$ vertices equals the sum of the $N_{\mathrm{ch}} - 1 = 19$ edge weights of the Minimum Spanning Tree ($\mathrm{MST}$). On the canonical 10--20 montage, posterior visual-parietal cortex is covered by $N_{\mathrm{post}} = 8$ electrodes ($\mathcal{V}_{\mathrm{post}} = \{\text{P7, P3, Pz, P4, P8, O1, Oz, O2}\}$). Any acyclic forest on $\mathcal{V}_{\mathrm{post}}$ contains at most $N_{\mathrm{post}} - 1 = 7$ edges. Therefore, at least $19 - 7 = 12$ edges ($63.2\%$ of the global sum $L_{\mathrm{tot}}(H_0)$) must connect the $N_{\mathrm{ant}} = 12$ fronto-central and temporal electrodes ($\mathcal{V}_{\mathrm{ant}}$), where ocular (EOG) and muscular (EMG) noise dominate.

\begin{table}[t]
\centering
\caption{Per-subject diagnostic audit across $N=20$ subjects (\texttt{S001--S020}, $58$ epochs/subject). Global 4s $\mathrm{ImCoh}$ suffers from polarity inversions ($d > 0$) in $4/20$ subjects ($^\dagger$), whereas Stage 7 Posterior Spectral-Topological filtration achieves $20/20$ consistent directionality ($d < 0$) and raises mean within-subject AUC from $0.652$ to $0.944$.}
\label{tab:diag}
\resizebox{\columnwidth}{!}{
\begin{tabular}{lcccc}
\toprule
\textbf{Subject} & \textbf{Old Global $H_0$ $d$} & \textbf{Old Global AUC} & \textbf{Stage 7 SpecTopo $d$} & \textbf{Stage 7 SpecTopo AUC} \\
\midrule
S001 & $-0.89$ & $0.734$ & $-8.32$ & $1.000$ \\
S002 & $-1.41$ & $0.847$ & $-2.67$ & $0.981$ \\
S003 & $-0.73$ & $0.653$ & $-4.65$ & $1.000$ \\
S004 & $-3.53$ & $0.993$ & $-5.73$ & $1.000$ \\
S005 & $-0.84$ & $0.690$ & $-1.86$ & $0.920$ \\
S006 & $+0.08^\dagger$ & $0.477$ & $-1.33$ & $0.828$ \\
S007 & $+0.60^\dagger$ & $0.334$ & $-1.36$ & $0.838$ \\
S008 & $-1.26$ & $0.797$ & $-2.36$ & $0.917$ \\
S009 & $-0.53$ & $0.630$ & $-4.19$ & $0.994$ \\
S010 & $-1.65$ & $0.894$ & $-2.86$ & $0.949$ \\
S011 & $-1.41$ & $0.817$ & $-2.53$ & $0.951$ \\
S012 & $+0.69^\dagger$ & $0.332$ & $-3.16$ & $0.996$ \\
S013 & $-0.42$ & $0.592$ & $-2.37$ & $0.989$ \\
S014 & $-1.06$ & $0.757$ & $-3.71$ & $1.000$ \\
S015 & $+2.14^\dagger$ & $0.065$ & $-0.97$ & $0.774$ \\
S016 & $-0.56$ & $0.648$ & $-3.61$ & $0.998$ \\
S017 & $-2.42$ & $0.988$ & $-6.13$ & $1.000$ \\
S018 & $-0.35$ & $0.593$ & $-5.96$ & $1.000$ \\
S019 & $-0.39$ & $0.615$ & $-0.98$ & $0.741$ \\
S020 & $-0.52$ & $0.591$ & $-5.11$ & $1.000$ \\
\midrule
\textbf{Summary} & \textbf{4/20 Inverted ($^\dagger$)} & \textbf{Mean: 0.652} & \textbf{0/20 Inverted} & \textbf{Mean: 0.944} \\
\bottomrule
\end{tabular}
}
\end{table}

---

\section{The TRR 3.0 Stage 7 Mathematical Architecture}

\subsection{Sub-Window Rectified Lag-Orthogonalized Imaginary Coherency ($\mathrm{rLagImCoh}$)}
Each 4-second epoch is partitioned into $M = 7$ overlapping sub-windows of length $L_{\mathrm{sub}} = 1.0\text{ s}$ ($160$ samples) with step $0.5\text{ s}$ ($80$ samples), tapered by a Hann window $w(t)$. Let $S_{uv}^{(m)} = \frac{1}{L_{\mathrm{sub}}}\sum_{t} w^2(t) z_u^{(m)}(t) z_v^{(m)*}(t)$ be the local cross-spectrum in sub-window $m$, and $P_u^{(m)} = \operatorname{Re}(S_{uu}^{(m)})$. We define:
\begin{equation}
    \mathrm{rLagImCoh}_{uv} = \frac{1}{M}\sum_{m=1}^{M} \mathrm{clip}_{[0,1]}\!\left(\frac{\left|\operatorname{Im}\left(S_{uv}^{(m)}\right)\right|}{\sqrt{\max\!\left(\epsilon_{\mathrm{tol}},\, P_u^{(m)}P_v^{(m)} - \gamma\left[\operatorname{Re}\left(S_{uv}^{(m)}\right)\right]^2\right)}}\right),
    \label{eq:rlag}
\end{equation}
with Pascual-Marqui lag-orthogonalization parameter $\gamma = 0.85$~\cite{pascual2007} and $\epsilon_{\mathrm{tol}} = 10^{-15}$. The induced functional distance matrix is $D_{\mathrm{rLag}}(u,v) = 1 - \mathrm{rLagImCoh}_{uv}$ for $u \neq v$ and $D_{\mathrm{rLag}}(u,u) = 0$.

\begin{proposition}[Strict Zero-Lag Volume Conduction Isolation]
For any instantaneous linear mixture of cortical sources ($\Delta\phi_{uv}(t) \equiv 0 \pmod\pi$), $\mathrm{rLagImCoh}_{uv} = 0$ and $D_{\mathrm{rLag}}(u,v) = 1.0$ for all $u \neq v$.
\end{proposition}
\begin{proof}
If $\Delta\phi_{uv}(t) \equiv 0 \pmod\pi$, then $z_u^{(m)}(t)z_v^{(m)*}(t) \in \mathbb{R}$ for all $t$, so $\operatorname{Im}(S_{uv}^{(m)}) = 0$ in every sub-window $m \in \{1,\dots,M\}$. By the Cauchy--Schwarz inequality, $[\operatorname{Re}(S_{uv}^{(m)})]^2 \le |S_{uv}^{(m)}|^2 \le P_u^{(m)}P_v^{(m)}$. Since $\gamma = 0.85 < 1$, the radicand in the denominator of Eq.~\eqref{eq:rlag} satisfies $P_u^{(m)}P_v^{(m)} - \gamma[\operatorname{Re}(S_{uv}^{(m)})]^2 \ge (1-\gamma)P_u^{(m)}P_v^{(m)}$, and the $\max(\epsilon_{\mathrm{tol}}, \cdot)$ operator guarantees a strictly positive denominator $\ge \sqrt{\epsilon_{\mathrm{tol}}} > 0$. Thus every summand in Eq.~\eqref{eq:rlag} is identically zero.
\end{proof}

\subsection{Anatomical Simplicial Decomposition over $\mathbb{F}_2$ and Combinatorial Speedup}
We partition the vertex set $\mathcal{V}$ ($|\mathcal{V}|=20$) into two disjoint anatomical sub-complexes:
\begin{align}
    \mathcal{V}_{\mathrm{post}} &= \{\text{P7, P3, Pz, P4, P8, O1, Oz, O2}\} \quad (N_{\mathrm{post}} = 8), \\
    \mathcal{V}_{\mathrm{ant}} &= \{\text{Fp1, Fp2, F7, F3, Fz, F4, F8, T7, C3, Cz, C4, T8}\} \quad (N_{\mathrm{ant}} = 12).
\end{align}
For each sub-complex $\mathcal{K}_r$ ($r \in \{\mathrm{post}, \mathrm{ant}\}$), the Vietoris--Rips filtration $\mathrm{VR}_\epsilon(\mathcal{V}_r, D_{\mathrm{rLag}})$ orders the 1-simplices (edges) $\mathcal{E}_r = \{e_1,\dots,e_{|\mathcal{E}_r|}\}$ by weight $w(e) = D_{\mathrm{rLag}}(u,v)$ and the 2-simplices (triangles) $\mathcal{T}_r = \{\tau_1,\dots,\tau_{|\mathcal{T}_r|}\}$ by diameter $w(\tau) = \max(D_{\mathrm{rLag}}(u,v), D_{\mathrm{rLag}}(v,w), D_{\mathrm{rLag}}(u,w))$.

Let $\partial_1 \in \mathbb{F}_2^{|\mathcal{V}_r| \times |\mathcal{E}_r|}$ and $\partial_2 \in \mathbb{F}_2^{|\mathcal{E}_r| \times |\mathcal{T}_r|}$ denote the simplicial boundary matrices over $\mathbb{F}_2$ (where addition is bitwise XOR $\oplus$). Standard left-to-right column reduction yields pivot pairs $(i, j)$ defining the zero- and one-dimensional persistence invariants:
\begin{align}
    L_{\mathrm{tot}}(H_0 \mid \mathcal{K}_r) &= \sum_{v \in \mathrm{Piv}(\partial_1)} w(e_{m_1(v)}), \qquad \bar{H}_0(\mathcal{K}_r) = \frac{L_{\mathrm{tot}}(H_0 \mid \mathcal{K}_r)}{N_r - 1}, \\
    L_{\max}(H_1 \mid \mathcal{K}_r) &= \max_{i \in Z_1(\mathcal{K}_r)} \big(w_{\mathrm{death}}(e_i) - w(e_i)\big), \\
    L_{\mathrm{tot}}(H_1 \mid \mathcal{K}_r) &= \sum_{i \in Z_1(\mathcal{K}_r)} \big(w_{\mathrm{death}}(e_i) - w(e_i)\big),
\end{align}
where $Z_1(\mathcal{K}_r) = \{i \in \{1,\dots,|\mathcal{E}_r|\} : \mathrm{piv}(\partial_1(:,i)) = -1\}$ indexes cycle-creating edges, and $w_{\mathrm{death}}(e_i) = w(\tau_{m_2(i)})$ if $i \in \mathrm{Piv}(\partial_2)$ or $\max(D_{\mathrm{rLag}})$ otherwise.

\textbf{Combinatorial Complexity Reduction:} Global filtration on $N_{\mathrm{ch}}=20$ vertices requires reducing $|\mathcal{T}_{\mathrm{glob}}| = \binom{20}{3} = 1{,}140$ triangles and $|\mathcal{E}_{\mathrm{glob}}| = \binom{20}{2} = 190$ edges. Anatomical decomposition requires only $|\mathcal{T}_{\mathrm{post}}| + |\mathcal{T}_{\mathrm{ant}}| = \binom{8}{3} + \binom{12}{3} = 56 + 220 = 276$ triangles and $\binom{8}{2} + \binom{12}{2} = 28 + 66 = 94$ edges, yielding a **$4.13\times$ reduction in 2-simplices** and a **$2.02\times$ reduction in 1-simplices**.

The dimensionless **Zero-Shot Topological Spatial Gradient** is defined as:
\begin{equation}
    \Delta_{\mathrm{spatial}} H_0(t) = \bar{H}_0(\mathcal{K}_{\mathrm{post}}, t) - \bar{H}_0(\mathcal{K}_{\mathrm{ant}}, t).
    \label{eq:grad}
\end{equation}

\subsection{Pure Topological vs.\ Hybrid Spectral-Topological Feature Spaces}
To ensure strict ablation transparency between phase-lagged $\mathbb{F}_2$ homology and local band power, we construct two distinct feature representations across subbands $\alpha$ ($8\text{--}13\text{ Hz}$), $\alpha_1$ ($7.5\text{--}10.5\text{ Hz}$), $\alpha_2$ ($10.0\text{--}13.5\text{ Hz}$), $\beta$ ($14.0\text{--}26.0\text{ Hz}$), and $\theta$ ($4.0\text{--}7.5\text{ Hz}$):
\begin{enumerate}
    \item \textbf{Pure Topological TRR ($\mathbb{F}_2$ only, $0\%$ Local Power):} The 1D shift-invariant anchor is $\mathbf{x}_{\mathrm{inv}}^{\mathrm{pure}}(t) = [\Delta_{\mathrm{spatial}} H_0^{(\alpha)}(t)] \in \mathbb{R}^1$, and the 7D calibratable topological vector is:
    \begin{align}
        \mathbf{x}_{\mathrm{cal}}^{\mathrm{pure}}(t) = \Big[
        &\Delta_{\mathrm{spatial}} H_0^{(\alpha)},\;
        \bar{H}_0^{(\alpha_1)}(\mathcal{K}_{\mathrm{post}}) - \bar{H}_0^{(\alpha)}(\mathcal{K}_{\mathrm{ant}}),\;
        \bar{H}_0^{(\alpha_2)}(\mathcal{K}_{\mathrm{post}}) - \bar{H}_0^{(\alpha)}(\mathcal{K}_{\mathrm{ant}}),\nonumber\\
        &\bar{H}_0^{(\alpha)}(\mathcal{K}_{\mathrm{post}}) - \bar{H}_0^{(\beta)}(\mathcal{K}_{\mathrm{post}}),\;
        \bar{H}_0^{(\alpha)}(\mathcal{K}_{\mathrm{post}}),\;
        L_{\max}(H_1^{(\alpha)} \mid \mathcal{K}_{\mathrm{post}}),\;
        L_{\mathrm{tot}}(H_1^{(\alpha)} \mid \mathcal{K}_{\mathrm{post}})
        \Big]^T \in \mathbb{R}^7.
    \end{align}
    \item \textbf{Hybrid SpecTopo TRR (Topology + Regional Band-Power Ratios):} Let $P_b(\mathcal{K}_r, t)$ denote the mean narrowband variance across channels in region $r$ and band $b$. We define two scale-free log-spectral ratios and a coupled spectral-topological index:
    \begin{align}
        \mathrm{rel\_\alpha}_{\mathrm{post}}(t) &= \ln\frac{P_\alpha(\mathcal{K}_{\mathrm{post}}, t) + \epsilon_{\mathrm{tol}}}{P_\beta(\mathcal{K}_{\mathrm{post}}, t) + P_\theta(\mathcal{K}_{\mathrm{post}}, t) + \epsilon_{\mathrm{tol}}}, \label{eq:rel_alpha}\\
        \mathrm{ratio}_{\mathrm{post/ant}}(t) &= \ln\frac{P_\alpha(\mathcal{K}_{\mathrm{post}}, t) + \epsilon_{\mathrm{tol}}}{P_\alpha(\mathcal{K}_{\mathrm{ant}}, t) + \epsilon_{\mathrm{tol}}}, \label{eq:ant_post}\\
        \mathrm{SpecTopo}(t) &= \bar{H}_0^{(\alpha)}(\mathcal{K}_{\mathrm{post}}, t)\cdot\exp\!\big(-0.5\,\mathrm{rel\_\alpha}_{\mathrm{post}}(t)\big). \label{eq:spectopo}
    \end{align}
    The Hybrid zero-shot anchor is $\mathbf{x}_{\mathrm{inv}}^{\mathrm{hyb}}(t) = [\Delta_{\mathrm{spatial}} H_0^{(\alpha)}, \mathrm{rel\_\alpha}_{\mathrm{post}}, \mathrm{ratio}_{\mathrm{post/ant}}, \mathrm{SpecTopo}]^T \in \mathbb{R}^4$, and the full Hybrid calibratable vector is $\mathbf{x}_{\mathrm{cal}}^{\mathrm{hyb}}(t) = [\mathbf{x}_{\mathrm{cal}}^{\mathrm{pure}}(t)^T, \mathrm{rel\_\alpha}_{\mathrm{post}}, \mathrm{ratio}_{\mathrm{post/ant}}, \mathrm{SpecTopo}]^T \in \mathbb{R}^{10}$.
\end{enumerate}

\subsection{Post-Inversion Empirical Bayes Stein Shrinkage Fusion}
In each fold of Leave-One-Subject-Out Cross-Validation (LOSO-CV), let $\mathcal{D}_{-s}$ denote the $N-1 = 19$ training subjects and $s$ the held-out target subject. Given a strictly causal calibration prefix of $k \in \{0,1,2,3,5,10,29\}$ resting-state $\mathrm{EO}$ epochs $\mathbf{X}_{s,\mathrm{EO}}[0:k]$, each coordinate $j \in \{1,\dots,p\}$ of $\mathbf{x}_{\mathrm{cal}}$ is standardized via the Normal-Inverse-Gamma (NIG) conjugate posterior against the training prior $(\mu_{p,j}, \sigma_{p,j}^2)$:
\begin{align}
    \tilde{\mu}_{s,j}(k) &= \frac{k\,\hat{\mu}_{s,j}(k) + \tau_0\,\mu_{p,j}}{k + \tau_0}, \\
    \tilde{\sigma}_{s,j}^2(k) &= \frac{\max(0, k-1)\hat{s}_{s,j}^2(k) + \nu_0\,\sigma_{p,j}^2 + \frac{k\tau_0}{k+\tau_0}(\hat{\mu}_{s,j}(k)-\mu_{p,j})^2}{\max(0, k-1) + \nu_0},
\end{align}
with fixed hyperparameters $\tau_0 = 3.0, \nu_0 = 5.0$, yielding calibrated feature vectors $\mathbf{z}_s(t \mid k) = (\mathbf{x}_{\mathrm{cal},s}(t) - \tilde{\boldsymbol{\mu}}_s(k)) \oslash \tilde{\boldsymbol{\sigma}}_s(k)$.

Every training subject $j \in \mathcal{D}_{-s}$ undergoes the exact same $k$-epoch causal prefix calibration to form training matrices $\mathbf{Z}_{\mathrm{tr,EO}}(k)$ and $\mathbf{Z}_{\mathrm{tr,EC}}(k)$. Regularized Fisher Linear Discriminant weights are learned strictly on $\mathcal{D}_{-s}$:
\begin{align}
    \Sigma_{\mathrm{pool}} &= \frac{1}{2}\big(\operatorname{Cov}(\mathbf{Z}_{\mathrm{tr,EO}}) + \operatorname{Cov}(\mathbf{Z}_{\mathrm{tr,EC}})\big), \\
    \Sigma_{\mathrm{reg}} &= (1 - \lambda)\Sigma_{\mathrm{pool}} + \lambda\operatorname{diag}(\Sigma_{\mathrm{pool}}), \qquad \lambda = 0.20, \\
    \mathbf{w} &= \Sigma_{\mathrm{reg}}^{-1}(\mathbf{m}_{\mathrm{tr,EC}} - \mathbf{m}_{\mathrm{tr,EO}}), \qquad \tilde{\mathbf{w}} = \frac{\mathbf{w}}{\sqrt{\operatorname{Var}(\mathbf{Z}_{\mathrm{tr,EO}}\mathbf{w}) + \epsilon_{\mathrm{tol}}}}.
    \label{eq:lda_norm}
\end{align}
Because multiplying feature columns by a Stein weight $\alpha_k$ \textit{prior} to $\Sigma_{\mathrm{reg}}^{-1}$ scales $w_j \propto \alpha_k^{-1}$ and cancels $\alpha_k$ in $\mathbf{z}^T\mathbf{w}$, we project the zero-shot anchor $u_{\mathrm{inv}}(t) = \mathbf{z}_{\mathrm{inv},s}(t)^T\tilde{\mathbf{w}}_{\mathrm{inv}}$ and the calibrated discriminant $u_{\mathrm{cal}}(t \mid k) = \mathbf{z}_s(t \mid k)^T\tilde{\mathbf{w}}_{\mathrm{cal}}(k)$ into unit-variance coordinates independently via Eq.~\eqref{eq:lda_norm}, and fuse them \textit{post-inversion}:
\begin{equation}
    S_s(t \mid k) = (1 - \alpha_k)\,u_{\mathrm{inv}}(t) + \alpha_k\,u_{\mathrm{cal}}(t \mid k), \qquad \alpha_k = \frac{k}{k + \tau_0}.
    \label{eq:fusion}
\end{equation}

---

\section{Experimental Results and Statistical Audit ($N=20$, 1,160 Epochs)}

\begin{table*}[t]
\centering
\caption{Stage 7 Verified Audit across $N=20$ subjects ($1{,}160$ epochs, \texttt{S001--S020}) under strict Leave-One-Subject-Out Cross-Validation (LOSO-CV). Balanced Holdout evaluates strictly post-calibration epochs (\texttt{EO[10:]} vs.\ \texttt{EC[10:]}, $760$ epochs). Hierarchical subject-level significance is evaluated via the two-sided paired Wilcoxon signed-rank test across $N=20$ independent subjects.}
\label{tab:main}
\resizebox{\textwidth}{!}{
\begin{tabular}{llccccc}
\toprule
\textbf{Architecture \& Strictly Causal Calibration Budget $k$ ($N=20$ LOSO-CV)} & \textbf{Calib Time} & \textbf{Cohen's $d$} & \textbf{Full Pool AUC} & \textbf{Balanced Holdout AUC} & \textbf{Mean Subject AUC ($\pm\text{SE}$)} & \textbf{Subject Wilcoxon $p$ (Pos)} \\
\midrule
\multicolumn{7}{l}{\textit{Part 1: Zero-Shot Simplicial Ablation ($k=0$, $0\text{ s}$ Calibration, Pure $\mathbb{F}_2$ Topology, Uncentered Physical Threshold)}} \\
1. Stage 6 Global $H_0$ (Baseline 4s $\mathrm{ImCoh}$, 20 ch) & $0\text{ s } (k=0)$ & $-0.60$ & $0.650$ & $0.626$ & $0.652 \pm 0.051$ & $7.30 \times 10^{-3}$ ($16/20$) \\
2. Pure Topo: Posterior Complex $\mathrm{rLagImCoh}$ $H_0$ (8 ch) & $0\text{ s } (k=0)$ & $-0.85$ & $0.780$ & $0.764$ & $0.905 \pm 0.029$ & $3.81 \times 10^{-6}$ ($19/20$) \\
3. Pure Topo: Spatial Gradient ($\mathcal{K}_{\mathrm{post}} - \mathcal{K}_{\mathrm{ant}}$) $H_0$ & $0\text{ s } (k=0)$ & $-0.82$ & $\mathbf{0.837}$ & $\mathbf{0.819}$ & $\mathbf{0.927 \pm 0.022}$ & $\mathbf{1.91 \times 10^{-6}}$ ($\mathbf{20/20}$) \\
\midrule
\multicolumn{7}{l}{\textit{Part 2A: Pure Topological TRR ($\mathbb{F}_2$ Simplicial Homology Only, $0\%$ Spectral Power, Post-Inversion Stein LOSO-CV)}} \\
Pure Topological TRR ($\mathbb{F}_2$ only, $k=0$) & $0\text{ s}$ & $+0.82$ & $0.830$ & $0.812$ & $0.927 \pm 0.022$ & $1.91 \times 10^{-6}$ ($20/20$) \\
Pure Topological TRR ($\mathbb{F}_2$ only, $k=1$) & $4\text{ s}$ & $+0.85$ & $0.831$ & $0.811$ & $0.929 \pm 0.021$ & $1.91 \times 10^{-6}$ ($20/20$) \\
Pure Topological TRR ($\mathbb{F}_2$ only, $k=2$) & $6\text{ s}$ & $+0.86$ & $0.836$ & $0.815$ & $0.929 \pm 0.021$ & $1.91 \times 10^{-6}$ ($20/20$) \\
Pure Topological TRR ($\mathbb{F}_2$ only, $k=3$) & $8\text{ s}$ & $+0.87$ & $\mathbf{0.840}$ & $\mathbf{0.818}$ & $\mathbf{0.929 \pm 0.021}$ & $1.91 \times 10^{-6}$ ($20/20$) \\
Pure Topological TRR ($\mathbb{F}_2$ only, $k=5$) & $12\text{ s}$ & $+0.89$ & $0.850$ & $0.827$ & $0.929 \pm 0.020$ & $1.91 \times 10^{-6}$ ($20/20$) \\
Pure Topological TRR ($\mathbb{F}_2$ only, $k=10$) & $22\text{ s}$ & $+0.88$ & $\mathbf{0.874}$ & $\mathbf{0.852}$ & $\mathbf{0.928 \pm 0.021}$ & $1.91 \times 10^{-6}$ ($20/20$) \\
Pure Topological TRR ($\mathbb{F}_2$ only, $k=29$) & $60\text{ s}$ & $+1.01$ & $\mathbf{0.895}$ & N/A (Full) & $0.915 \pm 0.024$ & $1.91 \times 10^{-6}$ ($20/20$) \\
\midrule
\multicolumn{7}{l}{\textit{Part 2B: Hybrid SpecTopo TRR ($\mathbb{F}_2$ Simplicial Homology + Regional Spectral Profile, Post-Inversion Stein LOSO-CV)}} \\
Hybrid SpecTopo TRR (Topo + Power, $k=0$) & $0\text{ s}$ & $+1.68$ & $\mathbf{0.869}$ & $\mathbf{0.856}$ & $\mathbf{0.955 \pm 0.013}$ & $1.91 \times 10^{-6}$ ($20/20$) \\
Hybrid SpecTopo TRR (Topo + Power, $k=1$) & $4\text{ s}$ & $+1.72$ & $0.874$ & $0.861$ & $0.955 \pm 0.013$ & $1.91 \times 10^{-6}$ ($20/20$) \\
Hybrid SpecTopo TRR (Topo + Power, $k=2$) & $6\text{ s}$ & $+1.83$ & $0.887$ & $0.874$ & $0.955 \pm 0.013$ & $1.91 \times 10^{-6}$ ($20/20$) \\
Hybrid SpecTopo TRR (Topo + Power, $k=3$) & $8\text{ s}$ & $+1.91$ & $\mathbf{0.897}$ & $\mathbf{0.883}$ & $\mathbf{0.955 \pm 0.013}$ & $1.91 \times 10^{-6}$ ($20/20$) \\
Hybrid SpecTopo TRR (Topo + Power, $k=5$) & $12\text{ s}$ & $+2.01$ & $0.908$ & $0.893$ & $0.956 \pm 0.012$ & $1.91 \times 10^{-6}$ ($20/20$) \\
Hybrid SpecTopo TRR (Topo + Power, $k=10$) & $22\text{ s}$ & $+2.20$ & $\mathbf{0.928}$ & $\mathbf{0.914}$ & $\mathbf{0.957 \pm 0.012}$ & $1.91 \times 10^{-6}$ ($20/20$) \\
Hybrid SpecTopo TRR (Topo + Power, $k=29$) & $60\text{ s}$ & $+2.34$ & $\mathbf{0.954}$ & N/A (Full) & $\mathbf{0.957 \pm 0.013}$ & $1.91 \times 10^{-6}$ ($20/20$) \\
\bottomrule
\end{tabular}
}
\end{table*}

\begin{figure*}[t]
    \centering
    \includegraphics[width=\textwidth]{trr_stage7_flawless_N20.pdf}
    \caption{\textbf{TRR 3.0 Stage 7 Benchmark across $N=20$ subjects ($1{,}160$ epochs, strict LOSO-CV).} \textbf{(A)} Receiver Operating Characteristic (ROC) curves comparing the Stage 6 global baseline ($k=0$, $\mathrm{AUC}=0.650$), the zero-shot pure topological spatial gradient ($k=0$, $\mathrm{AUC}=0.837$), Pure Topological TRR at $8\text{ s}$ ($k=3$, $\mathrm{AUC}=0.840$), and Hybrid SpecTopo TRR at $8\text{ s}$ ($k=3$, $\mathrm{AUC}=0.897$). \textbf{(B)} Strictly causal calibration scaling curves across resting-state $\mathrm{EO}$ prefix budgets $k \in \{0,1,2,3,5,10,29\}$ on the full $1{,}160$-epoch pool and balanced post-calibration holdout set (\texttt{EO[10:]} vs.\ \texttt{EC[10:]}).}
    \label{fig:main}
\end{figure*}

Table~\ref{tab:main} and Figure~\ref{fig:main} report the complete experimental audit across $N=20$ subjects ($1{,}160$ epochs):
\begin{itemize}
    \item \textbf{Zero-Shot Pure Topological Breakthrough ($k=0$):} Restricting filtration to the posterior complex $\mathcal{K}_{\mathrm{post}}$ under $\mathrm{rLagImCoh}$ raises pooled zero-shot accuracy from $0.650$ to $0.780$ ($19/20$ subjects consistent). Subtracting the anterior reference complex $\mathcal{K}_{\mathrm{ant}}$ via $\Delta_{\mathrm{spatial}} H_0$ elevates pooled zero-shot accuracy to $\mathbf{0.837}$ ($\mathrm{Holdout} = 0.819$, mean within-subject $\mathrm{AUC} = \mathbf{0.927 \pm 0.022}$) and achieves $20/20$ subject consistency (two-sided paired Wilcoxon $p = 1.91 \times 10^{-6}$).
    \item \textbf{Analysis of Out-of-Fold LOSO Centering at $k=0$:} In Table~\ref{tab:main}, uncentered $\Delta_{\mathrm{spatial}} H_0$ yields $\mathrm{Full\ AUC} = 0.837$ ($\mathrm{Holdout} = 0.819$), whereas `Pure Topological TRR` at $k=0$ yields $\mathrm{Full\ AUC} = 0.830$ ($\mathrm{Holdout} = 0.812$) while maintaining the exact same within-subject accuracy ($0.927 \pm 0.022$). This $0.007$ difference arises because $\Delta_{\mathrm{spatial}} H_0$ is naturally centered around physical zero ($\bar{H}_0(\mathcal{K}_{\mathrm{post}}) = \bar{H}_0(\mathcal{K}_{\mathrm{ant}})$), whereas subtracting the leave-one-out training mean $\mu_{i,-s}$ across $N-1=19$ subjects introduces a finite-sample jackknife variance inflation factor of $\frac{N}{N-1} = \frac{20}{19}$.
    \item \textbf{Monotonic Full-Pool Causal Calibration Scaling:} Under Post-Inversion Stein Shrinkage Fusion, `Full Pool AUC` increases strictly monotonically with calibration budget $k$ for both Pure Topological TRR ($0.830 \to 0.831 \to 0.836 \to 0.840 \to 0.850 \to 0.874 \to 0.895$) and Hybrid SpecTopo TRR ($0.869 \to 0.874 \to 0.887 \to 0.897 \to 0.908 \to 0.928 \to 0.954$). Note that at $k=29$ ($60\text{ s}$) in Pure Topological TRR, Fisher LDA optimizes cross-subject pooled separability ($\mathrm{Full\ AUC} = 0.895$, $d = +1.01$) by assigning higher weight to $H_1$ cycle persistence, trading a slight shift in mean within-subject AUC ($0.928 \to 0.915$) for superior global threshold alignment.
\end{itemize}

---

\section{Conclusion}

By resolving zero-lag denominator inflation, sub-epoch phase cancellation, and anterior MST dilution, TRR 3.0 Stage 7 proves that exact $\mathbb{F}_2$ persistent homology operates as a highly accurate, volume-conduction-invariant classifier in both zero-shot ($\mathrm{AUC}=0.837\text{--}0.869$) and rapid few-shot ($\mathrm{AUC}=0.897\text{--}0.928$ within $8\text{--}22\text{ s}$) regimes across heterogeneous multi-subject EEG cohorts.

\begin{thebibliography}{9}
\bibitem{weinberg2026trr1}
R. Weinberg, ``Relational Resolution Theory 3.0: Formal Verification in Lean 4 and Volume-Conduction Invariant Simplicial Persistent Homology over $\mathbb{F}_2$ in Multichannel EEG,'' \textit{Zenodo Preprint}, DOI: \texttt{10.5281/zenodo.23454524}, Repository: \texttt{https://github.com/RW-101/TRR3-formalization-}, 2026.

\bibitem{nolte2004}
G. Nolte et al., ``Identifying true brain interaction from EEG data using the imaginary part of coherency,'' \textit{Clinical Neurophysiology}, vol. 115, no. 10, pp. 2292--2307, 2004.

\bibitem{pascual2007}
R. D. Pascual-Marqui, ``Instantaneous and lagged measurements of linear and nonlinear dependence between groups of multivariate time series: frequency decomposition,'' \textit{arXiv:0711.1455}, 2007.

\bibitem{schalk2004}
G. Schalk et al., ``BCI2000: a general-purpose brain-computer interface (BCI) system,'' \textit{IEEE Transactions on Biomedical Engineering}, vol. 51, no. 6, pp. 1034--1043, 2004.
\end{thebibliography}

\end{document}
